\chapter{Conclusion}

\section{Summary of Contributions}

This project has delivered three interconnected contributions to the supervision and control
of the ICAM thermal process installation.

\paragraph{Real-time web supervision.}
A Docker-based Flask application communicates with the Siemens S7-1500 PLC via two
complementary protocols: Snap7 for direct data block access (DB1) and Siemens Open Pipe for
WinCC Unified tag exchange over a Unix socket. The application provides four web pages
accessible on the laboratory network: a full-screen P\&ID synoptique optimised for the
7\textquotedblright{} SIMATIC Unified touch panel, a cascade identification and commissioning
page for the engineer, a debug process tag monitor, and a historical trend chart with CSV
export. Socket.IO delivers real-time updates to all connected clients every 3\,s without
explicit polling.

\paragraph{Experimental process identification.}
An open-loop step-response identification experiment has been designed, implemented, and
conducted on the inner process ($F_{1,SP} \rightarrow T_{in1,HE1}$). The \code{StepIdentifier}
class automates the three-phase test (baseline averaging, step application, FOPTD analysis),
extracts the parameters $K$, $\tau$, $\theta$ by the 63.2\,\% method, and immediately
calculates IMC-tuned PI gains at three aggressiveness levels. The engineer can apply any
tuning level with a single click from the web interface.

\paragraph{Cascade PI+P controller.}
A cascade controller has been implemented in Python with: an inner PI loop in velocity
form with proportional action on measurement (bumpless transfer, implicit anti-windup);
an outer proportional controller driving the cold outlet temperature setpoint; and a
three-mode state machine (Manual / Auto-Inner / Auto-Full) with safe mode transitions.
The controller runs in a background thread and writes $F_{1,SP}$ to the PLC via Open Pipe
every 3\,s.

\section{Results Assessment}

\pendingbox{%
  To be completed after experimental sessions: identification results ($K$, $\tau$, $\theta$),
  closed-loop step response metrics, and cascade steady-state performance.%
}

The web supervision interface has been verified operationally: the Synoptique renders
correctly in the panel browser, Socket.IO connectivity is established, and the identification
workflow has been validated in simulation mode. The cascade controller mode state machine,
bumpless transfer, and anti-windup mechanisms have been verified through software testing.

\section{Limitations}
\label{sec:limitations}

\begin{description}[leftmargin=2cm, style=nextline, itemsep=6pt]

  \item[3-second sampling period]
    The Open Pipe poll rate of \SI{3}{\second} limits identification resolution to
    $\pm\SI{3}{\second}$ on both $\tau$ and $\theta$. For processes with dead times
    shorter than $\sim\SI{6}{\second}$, this may be insufficient for accurate IMC tuning.
    For the current thermal process ($\theta \gg \SI{6}{\second}$ expected), the limitation
    is not critical.

  \item[P outer loop offset]
    The outer proportional controller introduces a non-zero steady-state error in
    $T_{out2,HE1}$ (Section~6.3.2). Precise regulation requires either manual trimming of
    $T_{in1,SP}^{\text{base}}$ or an upgrade to a PI outer loop.

  \item[Open Pipe validation]
    The Open Pipe communication is implemented based on the documented JSON protocol and
    tested in simulation mode. Full end-to-end validation with the WinCC Unified runtime
    on the panel is pending Industrial Edge licence acquisition.

  \item[No PLC-side safety interlocks]
    The web application does not implement hardware-level safety interlocks (high-temperature
    cutoff, flow-loss detection). Process safety relies entirely on the TIA Portal program.
    This is acceptable for a laboratory demonstrator but would need to be addressed for
    any production deployment.

\end{description}

\section{Perspectives}
\label{sec:perspectives}

\begin{description}[leftmargin=2cm, style=nextline, itemsep=6pt]

  \item[Outer loop PI upgrade]
    Once the inner loop is validated, a step-response identification experiment on the
    outer process ($T_{in1} \rightarrow T_{out2}$) would yield $K_2$, $\tau_2$, $\theta_2$.
    These parameters enable IMC tuning of a PI outer controller, eliminating the
    steady-state $T_{out2}$ offset without manual trimming.

  \item[Faster sampling]
    Reducing the Open Pipe poll rate from 3\,s to 1\,s (by offloading controller
    computation to a separate thread) would improve identification resolution and increase
    the achievable inner loop bandwidth.

  \item[Alarm management]
    Adding configurable alarm thresholds (high $T_{hin}$, low $F_1$, valve fault detection)
    with Socket.IO broadcast events would improve operator safety awareness.

  \item[Complete Industrial Edge deployment]
    Obtaining Industrial Edge licence credentials and completing the panel deployment would
    realise the original vision of an embedded supervision interface, independent of the
    development PC.

  \item[Data analytics]
    The append-only CSV log provides raw data but no analysis. A lightweight analytics
    endpoint (trend statistics, energy balance, identification quality metrics) would
    increase the value of the logged data.

\end{description}
